\section[Hauville]{L’indexation Hauville: adapter le langage RAMEAU au droit}

La singularité de la bibliothèque Cujas réside en grande partie dans l'adoption d'un système d'indexation tout à fait spécifique, peu utilisé et largement méconnu des autres établissements documentaires. La bibliothèque numérique s'appuie en effet sur l'indexation Hauville, une méthodologie originale qui doit son nom à sa conceptrice, Frédérique Hauville, qui a activement collaboré au projet de préfiguration de la première bibliothèque numérique de Cujas. 

Ce choix méthodologique s'explique par le fait qu'à sa créaltion la bibliothèque numérique de Cujas reposait sur des \enquote{mots-clés libres qui généraient des incohérences} documentaires. Pour y remédier, il n'était cependant pas toujours envisagealbe de recourir au langage standardisé RAMEAU, considéré comme inadapté au droit en raison de caractère trop généraliste. C'est dans ce contexte que l'indexation Hauville a été développée : s'appuyant strucutrellement sur le référentiel RAMEAU, elle ne le mobilise que lorsqu'il s'avère pertinent, mais n'hésite pas à le déconstruire dans el cas contraire afin de proposer des vedettes plus adaptées. Bien que ce dispositif ait initialement été  validé par la BNF, cette approbation s'est doublée d'une mise en garde prospective qu'à sa pérennité. La BNF soulignait en effet qu'un tel système d'autorités commencerait inévitablement à poser de sérieux problèmes de cohérence et de gestion documentaire dès lors que la bibliothèque numérique franchirait le seuil critique des 1 000 documents numérisés.

L'indexation Hauville débute par une phase d'analyse documentaire approfondie, au cours de laquelle le catalogueur examine le titre, la table des matière, et la quatrième de couverture, tout en procédant à l'identification des thèmes de l'ouvrage. Une fois ces concepts sémantiques identifiés, ils sont traduits dans le langage d'autorité national RAMEAU. En l'absence de terme adéquat au sein de ce référentiel, l'établissement s'appuie \enquote{sur une forme normalisée qui est la plus courante comme Wikipédia ou vocabulaire des termes juridiques de Cornu par exemple} afin de combler les lacunes terminologiques du dictionnaire officiel. 

A l'issue de cette phase de normalisation, le technicien construit une védette d'indexation structurée en deux ou trois éléments en fonction de la spécialité du domaine traité. Puis, elle construit la vedette pouvant être composé entre deux ou trois éléments en fonction du domaine traité. L'indexation Hauville comprend la vedette de forme (code, travaux préparatoires...), et le concept juridique. L'idée est de pouvoir \enquote{relire la vedette à l'envers en mettant un de ou dans le(s)}. Un exemple type de facette pouvant être trouvé sur CujasNum est celle-ci : Eglise -- Administration. Une autre particularité de cette indexation est possibilité d'avoir \enquote{une indexation à la fois générique et spécifique. Il est possible d'avoir plusieurs vedettes construites ayant des liens de hiérarchie entre elles}, comme par exemple : Droit civil -- Obligations et Droit civil -- Contrats. Cela permet ainsi de mieux guider le lecteur.

À l'origine, l'application pratique de cette méthode d'indexation souffre d'un manque d'encadrement normatif au sein de la bibliothèque Cujas. Il était en effet fréquent qu'elle soit faite par différente personnes, car il s'agissait initialement\enquote{de mots-clés apposés dans les notices Dublin Core}. La multiplication des personnes en charge de l'indexation ont rapidement entraîné\enquote{ainsi des incohérences} sémantiques majeures. Certains catalogueurs utilisaient par exemple un simple trait d'union tandis d'autres recouret au tiret cadratin, une divergence typographique en apparence mineure mais suffisante pour générer artificiellement de nouvelles facettes redondantes et non connectées dans l'indexeur XTF. Pour pallier cette dérive et homogénéiser le catalogue, la direction impose un cadre strict en exigeant que \enquote{tout le monde indexerait de la même façon, en proposant de recourir à la liste d'autorité RAMEAU sur BNOpale}.Face à ces difficultés structurelles, la BNF recommande de \enquote{faire une liste fermée du vocabulaire utilisé RAMEAU tant qu'il n'y a pas trop d'ouvrages numérisés}. La BNF émet toutefois un avertissement quant à la viabilité à long terme de cette stratégie de contournement, car \enquote{à partir de milliers de documents, cela n'est plus possibile}. 

Au sein d'une bibliothèque numérique, l'indexation s'impose comme un élément pivot d'une navigation intuitive et efficace. Elle émane directement de la qualité des métadonnées, lesquelles demeurent les \enquote{garants de fonctionnalités de recherche et de navigation satisfaisantes dans une bibliothèque numérique}. Dès lors, l'interopérabilité des données s'avère cruciale : privée de cette capacité d'échange standardisé, une bibliothèque numérique risque son évincement progressif de la communauté malgré tous les services novateurs que celle-ci peut proposer. C'est précisément à l'aune de ce double impératif de description et d'interconnexion qu'il convient d'analyser les défaillances de l'infrastructure de CujasNum, dont l'obsolescence technique a progressivement transformé l'innovation initiale en une lourde dette technique.